% Part: first-order-logic
% Chapter: natural-deduction
% Section: proving-things-quant

\documentclass[../../../include/open-logic-section]{subfiles}

\begin{document}

\olfileid{fol}{ntd}{prq}

\olsection{పరిమాణీకరణికాలతో \usetoken{P}{derivation}}

\begin{ex}
పరిమాణీకరణికాలతో వ్యవహరించేటప్పుడు ఐగెన్ చరరాశి షరతును ఉల్లంఘించకుండా
చూసుకోవాలి. అందుకోసం కొన్నిసార్లు కొన్ని నిగమనాలను చేసే క్రమాన్ని
మార్చి ప్రయత్నించాలి. సాధారణంగా ఐగెన్ చరరాశి షరతుకు లోబడే నియమాలను
ముందుగా చూసుకోవడం ఉపయోగకరం (పూర్తయిన నిరూపణలో అవి కిందివైపు ఉంటాయి).

$\lexists[x][\lnot !A(x)] \lif \lnot \lforall[x][!A(x)]$ అనే
!!{formula}కు ఒక !!a{derivation} ఎలా ఇవ్వాలో చూద్దాం. ఎప్పటిలాగే ఇలా
రాయడంతో మొదలుపెడతాం:
\begin{prooftree}
\AxiomC{}
\UnaryInfC{$\lexists[x][\lnot !A(x)]\lif \lnot \lforall[x][!A(x)]$}
\end{prooftree}
చివరి మెట్టును \Intro{\lif} నియమంతో సమర్థించడానికి ఏమి కావాలో మొదట
రాసుకుందాం.
\begin{prooftree}
\AxiomC{$\Discharge{\lexists[x][\lnot !A(x)]}{1}$}
\DeduceC{$\lnot \lforall[x][!A(x)]$}
\DischargeRule{\Intro{\lif}}{1}
\UnaryInfC{$\lexists[x][\lnot !A(x)]\lif \lnot \lforall[x][!A(x)]$}
\end{prooftree}
$\lnot \lforall[x][!A(x)]$కు వర్తింపజేయదగిన స్పష్టమైన నియమం లేదు
కనుక \Elim{\lexists} నియమాన్ని వాడగలిగేలా !!{derivation}ను
అమర్చుకుందాం. ఇక్కడ ఐగెన్ చరరాశి షరతును గమనించి,
$\lexists[x][!A(x)]$లో గానీ అది ఆధారపడే పరికల్పనల్లో గానీ కనిపించని
స్థిరాంకాన్ని ఎంచుకోవాలి. (అయితే ఏ !!{constant}s కనిపించడం లేదు కనుక
ఏ ఎంపికైనా సరిపోతుంది.)
\begin{prooftree}
\AxiomC{$\Discharge{\lexists[x][\lnot !A(x)]}{1}$}
\AxiomC{$\Discharge{\lnot !A(a)}{2}$}
\DeduceC{$\lnot \lforall[x][!A(x)]$}
\DischargeRule{\Elim{\lexists}}{2}
\BinaryInfC{$\lnot \lforall[x][!A(x)]$}
\DischargeRule{\Intro{\lif}}{1}
\UnaryInfC{$\lexists[x][\lnot !A(x)] \lif \lnot \lforall[x][!A(x)]$}
\end{prooftree}
$\lnot \lforall[x][!A(x)]$ను నిగమించడానికి \Intro{\lnot} నియమాన్ని
వాడటానికి ప్రయత్నిద్దాం: ఇందుకోసం $\lforall[x][!A(x)]$ను అదనపు
పరికల్పనగా వాడుకొని అయినా ఒక వైరుధ్యాన్ని నిగమించాలి. ఈ వైరుధ్యంలో
\Elim{\lexists} నిగమనం ఉపసంహరించే $\lnot !A(a)$ పరికల్పన కూడా
ఉండవచ్చు. దాన్ని ఇలా అమర్చవచ్చు:
\begin{prooftree}
\AxiomC{$\Discharge{\lexists[x][\lnot !A(x)]}{1}$}
\AxiomC{$\Discharge{\lnot !A(a)}{2}, \Discharge{\lforall[x][!A(x)]}{3}$}
\DeduceC{$\lfalse$}
\DischargeRule{\Intro{\lnot}}{3}
\UnaryInfC{$\lnot \lforall[x][!A(x)]$}
\DischargeRule{\Elim{\lexists}}{2}
\BinaryInfC{$\lnot \lforall[x][!A(x)]$}
\DischargeRule{\Intro{\lif}}{1}
\UnaryInfC{$\lexists[x][\lnot !A(x)]\lif \lnot \lforall[x][!A(x)]$}
\end{prooftree}
వైరుధ్యాన్ని పొందడానికి దగ్గరపడ్డాం. ఐగెన్ చరరాశి షరతులేమీ లేని
\Elim{\lforall}ను వర్తింపజేయడం అత్యంత సులభం. సార్వత్రికంగా
పరిమాణీకరించిన $x$ స్థానంలో మనకు కావలసిన ఏ పదాన్నైనా వాడవచ్చు కనుక,
వైరుధ్యానికి చేరడానికి $a$నే కొనసాగించి వాడడం సముచితం.
\begin{prooftree}
\AxiomC{$\Discharge{\lexists[x][\lnot !A(x)]}{1}$}
\AxiomC{$\Discharge{\lnot !A(a)}{2}$}
\AxiomC{$\Discharge{\lforall[x][!A(x)]}{3}$}
\RightLabel{\Elim{\lforall}}
\UnaryInfC{$!A(a)$}
\RightLabel{\Elim{\lnot}}
\BinaryInfC{$\lfalse$}
\DischargeRule{\Intro{\lnot}}{3}
\UnaryInfC{$\lnot \lforall[x][!A(x)]$}
\DischargeRule{\Elim{\lexists}}{2}
\BinaryInfC{$\lnot \lforall[x][!A(x)]$}
\DischargeRule{\Intro{\lif}}{1}
\UnaryInfC{$\lexists[x][\lnot !A(x)]\lif \lnot \lforall[x][!A(x)]$}
\end{prooftree}

ముఖ్యంగా పరిమాణీకరణికాలతో వ్యవహరిస్తున్నప్పుడు, ఐగెన్ చరరాశి షరతు
ఉల్లంఘించబడలేదని ఈ దశలో మళ్లీ పరిశీలించడం ముఖ్యం. మనం వర్తింపజేసిన
ఐగెన్ చరరాశి షరతుకు లోబడే ఏకైక నియమం \Elim{\exists}; ఐగెన్
చరరాశి~$a$ అది ఆధారపడే ఏ పరికల్పనలోనూ కనిపించదు. కనుక ఇది ఒక సరైన
!!{derivation}.
\end{ex}

\begin{ex}
కొన్నిసార్లు కొన్ని !!{formula}s నుంచి ఒక !!a{formula}ను నిగమించవచ్చు.
ఇలాంటి సందర్భాల్లో ఉపసంహరించని పరికల్పనలు ఉండవచ్చు. మన అంతిమ లక్ష్యంతో
పాటు పరికల్పనలను కూడా గుర్తుంచుకోవడం ముఖ్యం.

$\lexists[x][(!A(x) \land !B(x))]$, $\lforall[x][(!B(x) \lif
!C(x,b))]$ అనే పరికల్పనల నుంచి !!{formula}
$\lexists[x][!C(x,b)]$కు ఒక !!a{derivation} ఎలా ఇవ్వాలో చూద్దాం.
ఎప్పటిలాగే నిష్కర్షను అడుగున రాస్తాం.
\begin{prooftree}
\AxiomC{}
\UnaryInfC{$\lexists[x][!C(x,b)]$}
\end{prooftree}

మనకు రెండు పూర్వాధారాలున్నాయి. మొదటిదాన్ని వాడడానికి, అంటే
$\lexists[x][(!A(x) \land !B(x))]$ నుంచి
$\lexists[x][!C(x, b)]$ యొక్క ఒక !!a{derivation}ను కనుగొనడానికి,
\Elim{\lexists} నియమాన్ని వాడాలి. దానికి ఐగెన్ చరరాశి షరతు ఉంది కనుక
ఆ నియమాన్ని ముందుగా వర్తింపజేద్దాం. కిందిది వస్తుంది:
\begin{prooftree}
\AxiomC{$\lexists[x][(!A(x) \land !B(x))]$}
\AxiomC{$\Discharge{!A(a) \land !B(a)}{1}$}
\DeduceC{$\lexists[x][!C(x,b)]$}
\DischargeRule{\Elim{\lexists}}{1}
\BinaryInfC{$\lexists[x][!C(x,b)]$}
\end{prooftree}
మనం ఉపయోగిస్తున్న రెండు పరికల్పనల్లోనూ $!B$ ఉంది. ఈ దశలో
$!B(a)$ను విడదీయడానికి \Elim{\land}ను వర్తింపజేయడం ఉపయోగకరం.
\begin{prooftree}
\AxiomC{$\lexists[x][(!A(x) \land !B(x)])$}
\AxiomC{$\Discharge{!A(a)
\land !B(a)}{1}$}
\RightLabel{\Elim{\land}}
\UnaryInfC{$!B(a)$}
\DeduceC{$\lexists[x][!C(x,b)]$}
\DischargeRule{\Elim{\lexists}}{1}
\BinaryInfC{$\lexists[x][!C(x,b)]$}
\end{prooftree}

మనకు ఉపయోగపడే రెండవ పరికల్పన $\lforall[x][(!B(x) \lif
!C(x,b))]$. ఐగెన్ చరరాశి షరతు లేదు కనుక $x$ స్థానంలో
!!{constant}~$a$ను పెట్టడానికి \Elim{\lforall}ను వాడి
$!B(a) \lif !C(a, b)$ను పొందవచ్చు. ఇప్పుడు $!B(a) \lif !C(a,b)$,
$!B(a)$ రెండూ ఉన్నాయి. తరువాత \Elim{\lif} నియమాన్ని నేరుగా
వర్తింపజేయాలి.
\begin{prooftree}
\AxiomC{$\lexists[x][(!A(x) \land !B(x))]$}
\AxiomC{$\lforall[x][(!B(x) \lif !C(x,b))]$}
\RightLabel{\Elim{\lforall}}
\UnaryInfC{$!B(a) \lif !C(a,b)$}
\AxiomC{$\Discharge{!A(a)
\land !B(a)}{1}$}
\RightLabel{\Elim{\land}}
\UnaryInfC{$!B(a)$}
\RightLabel{\Elim{\lif}}
\BinaryInfC{$!C(a,b)$}
\DeduceC{$\lexists[x][!C(x,b)]$}
\DischargeRule{\Elim{\lexists}}{1}
\insertBetweenHyps{\hspace{-5em}}
\BinaryInfC{$\lexists[x][!C(x,b)]$}
\end{prooftree}
దాదాపు చేరుకున్నాం!{} \Intro{\lexists}ను ఒక్కసారి వర్తింపజేస్తే మన
లక్ష్యానికి చేరతాం.
\begin{prooftree}
\AxiomC{$\lexists[x][(!A(x) \land !B(x))]$}
\AxiomC{$\lforall[x][(!B(x) \lif !C(x,b))]$}
\RightLabel{\Elim{\lforall}}
\UnaryInfC{$!B(a) \lif !C(a,b)$}
\AxiomC{$\Discharge{!A(a)
\land !B(a)}{1}$}
\RightLabel{\Elim{\land}}
\UnaryInfC{$!B(a)$}
\RightLabel{\Elim{\lif}}
\BinaryInfC{$!C(a,b)$}
\RightLabel{\Intro{\lexists}}
\UnaryInfC{$\lexists[x][!C(x,b)]$}
\DischargeRule{\Elim{\lexists}}{1}
\insertBetweenHyps{\hspace{-5em}}
\BinaryInfC{$\lexists[x][!C(x,b)]$}
\end{prooftree}
ప్రతి మెట్టులోనూ ఐగెన్ చరరాశి షరతులను ఉల్లంఘించకుండా చూసుకున్నాం
కనుక ఇది ఒక సరైన !!{derivation} అని నిశ్చయించవచ్చు.
\end{ex}

\begin{ex}
$\lforall[x][!A(x)] \lif \lexists[y][!B(y)]$,
$\lnot\lexists[y][!B(y)]$ అనే పరికల్పనల నుంచి !!{formula}
$\lnot\lforall[x][!A(x)]$కు ఒక !!a{derivation} ఇవ్వండి. ఎప్పటిలాగే
లక్ష్య !!{formula}ను అడుగున రాస్తాం.
\begin{prooftree}
\AxiomC{}
\UnaryInfC{$\lnot\lforall[x][!A(x)]$}
\end{prooftree}
!!{derivation}లో చివరి వరుస నిషేధం కనుక \Intro{\lnot}ను
ప్రయత్నిద్దాం. ఇందుకోసం ఒక వైరుధ్యాన్ని ఎలా !!{derive} చేయాలో
కనుగొనాలి.
\begin{prooftree}
\AxiomC{$\Discharge{\lforall[x][!A(x)]}{1}$}
\DeduceC{$\lfalse$}
\DischargeRule{\Intro{\lnot}}{1}
\UnaryInfC{$\lnot\lforall[x][!A(x)]$}
\end{prooftree}
ఇంతవరకు బాగుంది. \Elim{\lforall}ను వాడవచ్చు; కానీ అది లక్ష్యాన్ని
పొందడంలో సహాయపడుతుందో స్పష్టం కాదు. బదులుగా మన పరికల్పనల్లో ఒకదాన్ని
వాడదాం. $\lforall[x][!A(x)] \lif \lexists[y][!B(y)]$తో పాటు
$\lforall[x][!A(x)]$ ఉంటే \Elim{\lif} నియమాన్ని వాడవచ్చు.
\begin{prooftree}
\AxiomC{$\lforall[x][!A(x)] \lif \lexists[y][!B(y)]$}
\AxiomC{$\Discharge{\lforall[x][!A(x)]}{1}$}
\RightLabel{\Elim{\lif}}
\BinaryInfC{$\lexists[y][!B(y)]$}
\DeduceC{$\lfalse$}
\DischargeRule{\Intro{\lnot}}{1}
\UnaryInfC{$\lnot\lforall[x][!A(x)]$}
\end{prooftree}
ఇప్పుడు వాడుకోవడానికి చివరి పరికల్పన ఒకటి మిగిలింది. \Elim{\lnot}ను
వాడి వైరుధ్యానికి చేరడానికి అది సహాయపడుతుందని కనిపిస్తోంది.
\begin{prooftree}
\AxiomC{$\lnot\lexists[y][!B(y)]$}
\AxiomC{$\lforall[x][!A(x)] \lif \lexists[y][!B(y)]$}
\AxiomC{$\Discharge{\lforall[x][!A(x)]}{1}$}
\RightLabel{\Elim{\lif}}
\BinaryInfC{$\lexists[y][!B(y)]$}
\RightLabel{\Elim{\lnot}}
\BinaryInfC{$\lfalse$}
\DischargeRule{\Intro{\lnot}}{1}
\UnaryInfC{$\lnot\lforall[x][!A(x)]$}
\end{prooftree}
\end{ex}

\begin{prob}
కిందివాటిని చూపే !!{derivation}s ఇవ్వండి:
\begin{enumerate}
\item $\Proves (\lforall[x][!A(x)] \land \lforall[y][!B(y)]) \lif
\lforall[z][(!A(z) \land !B(z))]$.
\item $\Proves (\lexists[x][!A(x)] \lor \lexists[y][!B(y)]) \lif
\lexists[z][(!A(z) \lor !B(z))]$.
\item $\lforall[x][(!A(x) \lif !B)] \Proves \lexists[y][!A(y)] \lif !B$.
\item $\lforall[x][\lnot !A(x)] \Proves \lnot\lexists[x][!A(x)]$.
\item $\Proves \lnot\lexists[x][!A(x)] \lif \lforall[x][\lnot !A(x)]$.
\item $\Proves \lnot\lexists[x][\lforall[y][((!A(x,y) \lif \lnot
!A(y,y)) \land (\lnot !A(y,y) \lif !A(x,y)))]]$.
\end{enumerate}
\end{prob}

\begin{prob}
కిందివాటిని చూపే !!{derivation}s ఇవ్వండి:
\begin{enumerate}
\item $\Proves \lnot\lforall[x][!A(x)] \lif \lexists[x][\lnot!A(x)]$.
\item $(\lforall[x][!A(x)] \lif !B) \Proves \lexists[y][(!A(y) \lif !B)]$.
\item $\Proves \lexists[x][(!A(x) \lif \lforall[y][!A(y)])]$.
\end{enumerate}
(వీటన్నింటికీ $\FalseCl$~నియమం అవసరం.)
\end{prob}

\end{document}
